% Generado por src/make_tables.py — no editar a mano.
\begin{table}[p]
\centering
\caption{Induction and recovery times (minutes after intraperitoneal injection of propofol \SI{7}{\milli\gram\per\kilo\gram} and ketamine \SI{50}{\milli\gram\per\kilo\gram}) in 13 male guinea pigs.}
\label{tab:times}
\small
\begin{tabular}{lcccc}
\toprule
Event & $n$ & Median (IQR) & Range & 95\% CI of median \\
\midrule
\multicolumn{5}{l}{\textit{Induction}} \\
\quad Reduced activity & 13 & 3 (2--4) & 2--11 & 2--5 \\
\quad Loss of righting reflex (LORR) & 13 & 5 (4--8) & 3--14 & 3--8 \\
\quad Loss of response to stimulus & 13 & 7 (4--11) & 3--14 & 4--12 \\
\quad Loss of pedal withdrawal reflex & 13 & 8 (4--10) & 3--14 & 4--11 \\
\quad Adequate anaesthetic plane & 12 & 10 (6--15) & 4--29 & 5--16 \\
\multicolumn{5}{l}{\textit{Recovery}} \\
\quad First movement & 13 & 48 (35--53) & 23--73 & 27--63 \\
\quad Return of pedal reflex & 13 & 56 (42--63) & 27--84 & 38--69 \\
\quad Return of righting reflex (RORR) & 13 & 61 (41--73) & 29--89 & 38--79 \\
\quad Sternal recumbency & 12 & 58 (40--72) & 31--97 & 38--83 \\
\quad Ambulation & 13 & 66 (46--93) & 35--147 & 42--93 \\
\quad Complete recovery & 13 & 108 (86--126) & 41--187 & 70--140 \\
\multicolumn{5}{l}{\textit{Duration}} \\
\quad Duration of LORR (RORR -- LORR) & 13 & 56 (33--67) & 23--81 & 24--76 \\
\quad Duration of pedal reflex absence & 13 & 52 (32--57) & 16--80 & 24--61 \\
\bottomrule
\end{tabular}
\par\smallskip\footnotesize IQR, interquartile range; CI, distribution-free confidence interval based on binomial order statistics.
\end{table}

\begin{table}[p]
\centering
\caption{Adverse events classified from the monitoring data with predefined thresholds (5--60 min after injection).}
\label{tab:adverse}
\small
\begin{tabular}{lccc}
\toprule
Event & $k/n$ & \% & 95\% CI (exact) \\
\midrule
Hypothermia (drop $\geq$1.0 \si{\celsius}) & 12/13 & 92 & 64--100 \\
Marked hypothermia (drop $\geq$2.0 \si{\celsius}) & 10/13 & 77 & 46--95 \\
Temperature $<$35.0 \si{\celsius} & 10/13 & 77 & 46--95 \\
Bradycardia ($\downarrow$$\geq$20 \% or $<$200 min$^{-1}$) & 4/13 & 31 & 9--61 \\
Tachycardia ($\uparrow$$\geq$20 \%) & 5/13 & 38 & 14--68 \\
Bradypnoea ($\downarrow$$\geq$50 \%) & 4/13 & 31 & 9--61 \\
Hypoxaemia (SpO$_2$ $<$90 \%)$^{\dagger}$ & 4/6 & 67 & 22--96 \\
Severe hypoxaemia (SpO$_2$ $<$80 \%)$^{\dagger}$ & 3/6 & 50 & 12--88 \\
\bottomrule
\end{tabular}
\par\smallskip\footnotesize Changes are relative to each animal's baseline. $^{\dagger}$Only animals with a physiologically plausible baseline SpO$_2$ ($\geq$90\%) while awake ($n=6$).
\end{table}
